\section{Analyse par NMF et K-means}

Deux méthodes non supervisées ont été utilisées pour explorer les régularités thématiques du corpus : la factorisation en matrices non négatives (NMF) et K-means. La première fait apparaître des topics à partir de groupes de mots fréquemment associés, tandis que la seconde rassemble les segments présentant des profils lexicaux proches. Leur comparaison permet d'observer le corpus selon deux perspectives complémentaires.

\subsection{La méthode NMF}

La NMF (\textit{Non-negative Matrix Factorization}) est une méthode de réduction dimensionnelle fréquemment utilisée pour l'exploration thématique. Elle décompose une matrice documents-termes en deux matrices non négatives : la première associe les segments aux topics, tandis que la seconde associe les topics aux termes qui les caractérisent.

Chaque segment peut ainsi être rattaché à plusieurs topics avec des poids différents. Cette représentation est adaptée aux textes littéraires, dans lesquels plusieurs thèmes peuvent apparaître simultanément au sein d'un même passage.

\subsubsection{Création de la matrice documents-termes}

Les segments lemmatisés ont d'abord été transformés en une représentation numérique à l'aide d'une pondération TF-IDF\footnote{Le TF-IDF (\textit{term frequency-inverse document frequency}) mesure l'importance d'un terme dans un document en tenant compte de sa fréquence dans l'ensemble du corpus. Le calcul a été réalisé avec la bibliothèque \texttt{scikit-learn}.}. Cette pondération valorise les termes caractéristiques d'un segment tout en réduisant l'influence des mots présents dans une grande partie du corpus.

Deux paramètres ont été utilisés pour filtrer le vocabulaire. Le paramètre \texttt{min\_df=3} écarte les termes présents dans moins de trois segments, considérés comme trop rares pour faire apparaître des régularités thématiques. À l'inverse, \texttt{max\_df=0.8} retire les termes présents dans plus de 80~\% des segments, qui sont trop fréquents pour être réellement discriminants\footnote{Le code du vectoriseur est disponible sur le dépôt GitHub associé au mémoire.}.

Plusieurs configurations ont été expérimentées. Les paramètres retenus offrent un compromis entre la réduction du vocabulaire et la conservation des termes utiles à l'interprétation. La matrice obtenue sert ensuite d'entrée au modèle NMF.

\subsubsection{Choix du nombre de topics}

Le nombre de topics influence directement la précision des résultats. Une valeur trop faible produit des catégories très générales, tandis qu'une valeur trop élevée fragmente les thèmes et multiplie les topics proches ou difficiles à interpréter.

La méthode du coude a d'abord été utilisée pour orienter ce choix\footnote{Le code correspondant est présenté en annexe.}. Plusieurs modèles ont été entraînés avec un nombre croissant de topics, puis comparés à partir de leur erreur de reconstruction. Un coude clairement marqué aurait indiqué qu'au-delà d'un certain seuil, l'ajout de nouveaux topics améliorait peu la reconstruction de la matrice.

\begin{figure}[H]
    \centering
    \includegraphics[width=0.8\textwidth]{images/chapitre_1/methodologie/coude_2_40.png}
    \caption{Méthode du coude pour le choix du nombre de topics dans la NMF, entre 2 et 40 topics}
    \label{fig:methode_du_coude_nmf_2_40}
\end{figure}

Comme le montre la figure~\ref{fig:methode_du_coude_nmf_2_40}, l'erreur diminue régulièrement lorsque le nombre de topics augmente. Cette évolution est attendue : un modèle disposant de davantage de topics peut reconstruire plus précisément la matrice initiale. La courbe ne présente cependant aucune rupture suffisamment nette pour désigner une valeur optimale.

Une seconde visualisation a donc été produite entre 20 et 39 topics afin d'examiner plus précisément cette partie de la courbe.

\begin{figure}[H]
    \centering
    \includegraphics[width=0.8\textwidth]{images/chapitre_1/methodologie/methode_du_coude_nmf.png}
    \caption{Méthode du coude pour le choix du nombre de topics dans la NMF, entre 20 et 39 topics}
    \label{fig:methode_du_coude_nmf_20_39}
\end{figure}

Ce second graphique confirme l'absence de coude clairement identifiable. Le choix final ne pouvait donc pas reposer uniquement sur l'erreur de reconstruction. Plusieurs modèles, comprenant notamment 10, 15, 20, 29, 30 et 35 topics, ont alors été comparés à partir des mots caractéristiques de chaque topic, de leur cohérence interne et de leur lisibilité.

Le modèle à 29 topics a finalement été retenu. Il offre un niveau de détail suffisant pour distinguer plusieurs ensembles thématiques sans produire une fragmentation excessive. Ce choix reste néanmoins empirique : il correspond au meilleur compromis observé dans le cadre de cette étude, et non à une valeur optimale démontrée de manière absolue.

\subsubsection{Paramétrage du modèle}

Le modèle final utilise le paramètre \texttt{n\_components=29}. Le paramètre \texttt{random\_state=42} assure la reproductibilité du traitement. L'initialisation \texttt{nndsvda}\footnote{L'initialisation \texttt{nndsvda} est une variante de NNDSVD dans laquelle les valeurs nulles sont remplacées par la moyenne de la matrice. Voir la documentation de \texttt{scikit-learn}.} fournit un point de départ stable, tandis que le solveur \texttt{cd}, pour \textit{coordinate descent}, optimise progressivement la factorisation.

La matrice \texttt{W} contient le poids de chacun des 29 topics dans chaque segment. La matrice \texttt{H} indique, pour sa part, le poids des termes dans chaque topic. L'examen des termes les mieux pondérés permet ensuite d'interpréter et de nommer les topics.

\subsubsection{Portée et limites de la NMF}

La NMF fournit une première cartographie des régularités lexicales du corpus. Elle permet notamment de faire apparaître des ensembles liés au travail, à la religion, à la famille, à l'argent, à la guerre, au commerce, à la médecine ou encore à l'espace domestique.

Les résultats restent toutefois dépendants des choix réalisés en amont : nettoyage, segmentation, lemmatisation, catégories grammaticales conservées, liste d'exclusion, paramètres du vectoriseur et nombre de topics. La modélisation thématique ne constitue donc pas une procédure entièrement automatique, mais une démarche exploratoire associant calcul statistique et interprétation humaine.

La nomination des topics comporte également une part de subjectivité. Deux lecteurs pourraient proposer des intitulés différents pour un même groupe de termes. Ces noms doivent par conséquent être compris comme des outils d'interprétation et non comme des catégories définitives.

Enfin, puisque le prétraitement conserve principalement les noms et les adjectifs, cette analyse porte avant tout sur le contenu thématique. Elle ne permet pas, à elle seule, de distinguer ce qui relève du naturalisme, de la signature personnelle de Zola ou des sujets particuliers de ses romans. Cette distinction sera étudiée à l'aide du corpus comparatif et de la classification automatique.

\subsection{La méthode K-means}

K-means a été utilisé comme une seconde manière d'explorer la structure du corpus. Contrairement à la NMF, qui représente chaque segment comme une combinaison de plusieurs topics, K-means attribue chaque segment à un seul cluster en fonction de sa proximité lexicale avec les autres segments.

L'objectif est de déterminer si les segments peuvent être regroupés en ensembles cohérents et si certains de ces regroupements présentent des ressemblances avec les topics obtenus par la NMF. Il ne s'agit pas d'une validation indépendante, puisque les deux méthodes utilisent la même représentation TF-IDF, mais d'une triangulation exploratoire fondée sur deux modes de structuration différents.

\subsubsection{Préparation de la matrice TF-IDF}

K-means est appliqué à la même matrice TF-IDF que la NMF. Les textes, la lemmatisation, la liste d'exclusion et les paramètres \texttt{min\_df} et \texttt{max\_df} restent donc identiques\footnote{Voir la section consacrée à la création de la matrice documents-termes.}. Ce choix garantit que les différences observées proviennent principalement du fonctionnement des deux méthodes et non d'un changement de représentation du corpus.

\subsubsection{Recherche du nombre de clusters}

Comme pour la NMF, le nombre de groupes doit être fixé avant l'entraînement du modèle. Une valeur trop faible risque de réunir des motifs distincts, tandis qu'une valeur trop élevée peut produire des clusters très proches ou difficiles à interpréter.

Plusieurs valeurs de \texttt{k}, comprises entre 2 et 40, ont été évaluées à l'aide du score de silhouette\footnote{Peter J. Rousseeuw, « Silhouettes: A Graphical Aid to the Interpretation and Validation of Cluster Analysis », \textit{Journal of Computational and Applied Mathematics}, vol.~20, 1987, p.~53--65.}. Ce score mesure à la fois la cohésion interne des clusters et leur séparation : une valeur élevée indique que les segments sont proches des éléments de leur propre groupe et éloignés des autres groupes.

\begin{figure}[H]
    \centering
    \includegraphics[width=0.8\textwidth]{images/chapitre_1/methodologie/silhouette.png}
    \caption{Évolution du score de silhouette en fonction du nombre de clusters dans K-means}
    \label{fig:silhouette_kmeans}
\end{figure}

La figure~\ref{fig:silhouette_kmeans} fournit une indication quantitative sur la structure du corpus, mais elle ne désigne pas à elle seule une valeur incontestable. Le modèle final a été fixé à 29 clusters afin de conserver un niveau de granularité comparable aux 29 topics de la NMF.

Ce choix facilite la confrontation des deux méthodes, mais il ne signifie pas qu'elles convergent indépendamment vers le nombre 29. Les 29 clusters ont été retenus en partie pour rendre la comparaison possible. L'analyse porte donc principalement sur les ressemblances et les différences entre les motifs lexicaux obtenus, et non sur une validation du nombre exact de groupes.

\subsubsection{Entraînement du modèle}

Le modèle final est entraîné avec \texttt{n\_clusters=29}. Le paramètre \texttt{random\_state=42} garantit la reproductibilité, tandis que \texttt{n\_init=10} lance l'algorithme avec plusieurs initialisations et conserve la meilleure solution.

À l'issue de l'entraînement, chaque segment est rattaché à l'un des 29 clusters. Les termes les plus caractéristiques de chaque groupe et les segments qui lui sont associés servent ensuite à interpréter son contenu.

\subsubsection{Portée et limites de K-means}

La principale limite de K-means tient à son fonctionnement exclusif : chaque segment appartient à un seul cluster. Or, un passage littéraire peut faire intervenir plusieurs thèmes simultanément, par exemple la famille, l'argent et la souffrance sociale. La NMF représente plus facilement ce type de coexistence, puisqu'elle attribue plusieurs poids thématiques à un même segment.

Les résultats de K-means sont également sensibles à la représentation TF-IDF et aux étapes de préparation du corpus. Une modification de la segmentation, de la lemmatisation ou du vocabulaire retenu pourrait produire une organisation différente. Comme pour la NMF, les clusters doivent donc être considérés comme une cartographie exploratoire plutôt que comme une classification définitive.

Enfin, leur nomination repose sur l'examen des termes et des segments les plus représentatifs. Cette interprétation humaine est nécessaire pour donner un sens littéraire aux groupes obtenus, mais elle introduit une part de subjectivité qui doit rester explicite.

\subsubsection{Bilan méthodologique}

La NMF et K-means proposent deux lectures complémentaires du corpus. La NMF décrit chaque segment comme une combinaison de topics, tandis que K-means construit une partition unique fondée sur la proximité lexicale. La comparaison de leurs résultats permet de repérer des motifs récurrents, mais ne constitue pas une validation indépendante puisque les deux méthodes reposent sur les mêmes données et la même matrice TF-IDF.

Une classification hiérarchique appliquée aux centres des clusters permet enfin de rassembler les 29 groupes en cinq familles thématiques plus générales. Cette organisation facilite la lecture des résultats, présentés dans la section suivante.